% from ega3-1-fr.tex
\oldpage[III]{152}

que $f_0$ est propre, on en conclut que $f_0$ est projectif (II, 5.5.4). Soit
\[
S=\bigoplus_{k\geq0}\mathfrak J^k/\mathfrak J^{k+1}
\]
la $A$-algèbre graduée associée à la filtration $\mathfrak J$-adique de $A$, qui est
de type fini puisque $A$ est noethérien~; si on pose
$\mathcal S'=f_0^*(\widetilde S)$,
$\mathcal M=\bigoplus_{k\geq0}\mathcal M_k$, $\mathcal S'$ est une
$\mathcal O_{X_0}$-Algèbre quasi-cohérente de type fini, et $\mathcal M$ un
$\mathcal S'$-Module gradué quasi-cohérent et de type fini (puisque $\mathcal F_0$
est cohérent et engendre le $\mathcal S'$-Module $\mathcal M$). On est donc dans les
conditions d'application du théorème (2.4.1, (ii)), et on en conclut qu'il existe
$n_0$ tel que, pour $n\geq n_0$ et pour tout $k$, on ait
\begin{equation}
\mathrm H^q(X_0,\mathcal M_k(n))=0\qquad\text{pour tout }q>0.
\tag{5.2.3.1}
\label{III.5.2.3.1-fr}
\end{equation}
On a donc aussi $\mathrm H^q(\mathfrak X,\mathcal M_k(n))=0$ pour $q>0$ et
$n\geq n_0$, $\mathcal M_k(n)$ étant cette fois considéré comme
$\mathcal O_{\mathfrak X}$-Module. Appliquant la suite exacte de cohomologie à
\[
0\to
\frac{\mathfrak J^k\mathcal F(n)}{\mathfrak J^{k+1}\mathcal F(n)}
\to
\frac{\mathfrak J^h\mathcal F(n)}{\mathfrak J^{k+1}\mathcal F(n)}
\to
\frac{\mathfrak J^h\mathcal F(n)}{\mathfrak J^k\mathcal F(n)}
\to0,
\]
on en déduit tout d'abord que pour $0\leq h<k$, $n\geq n_0$ et $q>0$, on a par
récurrence sur $k-h$
\begin{equation}
\mathrm H^q\left(\mathfrak X,
\frac{\mathfrak J^h\mathcal F(n)}{\mathfrak J^k\mathcal F(n)}\right)=0
\tag{5.2.3.2}
\label{III.5.2.3.2-fr}
\end{equation}
et en particulier pour $h=0$
\begin{equation}
\mathrm H^q(\mathfrak X,\mathcal F_k(n))=0
\qquad\text{pour }n\geq n_0,\ k\geq0\text{ et }q>0.
\tag{5.2.3.3}
\label{III.5.2.3.3-fr}
\end{equation}
Une autre portion de la suite exacte de cohomologie, pour $h=0$, donne la suite
exacte
\begin{equation}
\mathrm H^0(\mathfrak X,\mathcal F_{k+1}(n))\to
\mathrm H^0(\mathfrak X,\mathcal F_k(n))\to
\mathrm H^1\left(\mathfrak X,
\frac{\mathfrak J^k\mathcal F(n)}{\mathfrak J^{k+1}\mathcal F(n)}\right)=0
\tag{5.2.3.4}
\label{III.5.2.3.4-fr}
\end{equation}
d'où on déduit que pour $h\leq k$, l'application canonique
\begin{equation}
\mathrm H^0(\mathfrak X,\mathcal F_k(n))\longrightarrow
\mathrm H^0(\mathfrak X,\mathcal F_h(n))
\tag{5.2.3.5}
\label{III.5.2.3.5-fr}
\end{equation}
est surjective. Pour tout $q$, le système projectif
$(\mathrm H^q(\mathfrak X,\mathcal F_k(n)))_{k\geq0}$ vérifie donc la condition
(ML) pour $n\geq n_0$. Par ailleurs, tout ouvert formel affine $U$ de $\mathfrak X$
est aussi un ouvert affine dans chacun des $X_k$ (I, 10.5.2), donc on a
$\mathrm H^q(U,\mathcal F_k(n))=0$ pour tout $q>0$ (I, 1.3.1), et
$\mathrm H^0(U,\mathcal F_k(n))\to\mathrm H^0(U,\mathcal F_h(n))$ est surjective
pour $h\leq k$ (I, 1.3.9). Les conditions d'application de ($0$, 13.3.1) sont par
suite remplies, et on en conclut que, pour $n\geq n_0$~:
\begin{enumerate}
\item[$1^\circ$] Pour tout $q>0$,
\[
\mathrm H^q(\mathfrak X,\mathcal F(n))\longrightarrow
\varprojlim_k\mathrm H^q(\mathfrak X,\mathcal F_k(n))
\]
est bijectif, donc, en vertu de (\ref{III.5.2.3.3-fr}),
$\mathrm H^q(\mathfrak X,\mathcal F(n))=0$.
\item[$2^\circ$] L'homomorphisme
\[
\mathrm H^0(\mathfrak X,\mathcal F(n))\longrightarrow
\varprojlim_k\mathrm H^0(\mathfrak X,\mathcal F_k(n))
\]
est bijectif~; par ailleurs, comme les homomorphismes
(\ref{III.5.2.3.5-fr}) sont surjectifs, il en est de même de chacun des
homomorphismes
\[
\varprojlim_k\mathrm H^0(\mathfrak X,\mathcal F_k(n))
\longrightarrow\mathrm H^0(\mathfrak X,\mathcal F_h(n)),
\]
\end{enumerate}
ce qui achève la démonstration.
\end{proof}

\begin{corollary}[5.2.4]
\label{III.5.2.4-fr}
Les hypothèses étant celles de (\hyperref[III.5.2.3-fr]{5.2.3}), pour tout
$\mathcal O_{\mathfrak X}$-Module cohérent $\mathcal F$, il existe un entier $N$
tel que pour $n\geq N$, $\mathcal F(n)$ soit engendré par ses sections au-dessus

